\section{From noisy values to a bounded witness}\label{sec:search}
We prove Theorem~\ref{thm:main}. In this section $e=2^{-k}$.
For $v\ge1$, write $\clog(v)$ for the smallest $j\le v$ with $2^j\ge v$.
This is a bounded integer operation. Choose the proof precision
\[
 h=k+8+2\clog(d+1),\qquad \delta=2^{-h}
       \le\frac{e}{256d(d+1)},\qquad \tau=e/2 .
\]
It is finer than the output precision $t=k+6$.

\paragraph{Correcting each observed vector.}
For a rational $p\in P_\Q$ put $a_p=A(p,h)$ and define
\[
 v_p=a_p-\frac{p\cdot a_p}{p\cdot p}\,p .
\]
Since $(\sum |p_i|)^2\le d\sum p_i^2$, the denominator is at least
$1/d$. Dotting with $p$ gives exactly zero. Condition (T) and
$\len p=1$ give the explicit correction bound $\len{v_p-a_p}\le d\delta$.
At two points within $2^{-\mu(h)}$, condition (U) now gives
\[
 \len{v_p-v_q}\le(2d+3)\delta\le e/(32d)=\tau/(16d).
\]
This step only uses the individual errors and a same-precision
comparison. In particular, the correction never amplifies a large
field value in the estimate.

\paragraph{The fine witness used in the proof.}
Take $s=1+\max(\mu(h),\clog(8d))$ and $r=ns$. By
Lemma~\ref{lem:mesh} the mesh of $K_r$ is at most
$\min(2^{-\mu(h)},1/(8d))$. The preceding bounds verify the hypotheses
of Theorem~\ref{thm:finite-table} on its vertices. Thus the finite
table supplies a $p$ with $\len{v_p}\le e/2$, whence
$\len{A(p,h)}\le e/2+d\delta$.

\paragraph{Moving to the search grid.}
Write $p_i=\sigma_i b_i$ with $\sigma_i\in\{\pm1\}$, $b_i\ge0$,
$\sum b_i=1$ (choose $+1$ for zero). Set
\[
 c_i=\lfloor Db_i\rfloor\quad(i<d),\qquad
 c_d=D-\sum_{i<d}c_i,\qquad q_i=\sigma_i c_i/D .
\]
Then $c_d\ge0$ and $q\in P_\Q$. The discrepancy in the last
coordinate is the sum of the other rounding errors before restoring
signs. Hence $\len{p-q}\le2(d-1)/D\le2^{-\mu(t)}$.
Using (C) at $p$, then (U) between $p$ and $q$, gives
\[
 \len{A(q,t)}
 \le e/2+d\delta+2\cdot2^{-t}+3\cdot2^{-t}
 \le 3e/4 .
\]
For the last comparison, $d\delta\le e/256$ and $2^{-t}=e/64$.
This proves that at least one candidate passes the rational test.
The scan returns its first passing candidate in the chosen order.
For any such candidate and every $j\ge t$, (C) gives
$\len{A(q,j)}\le 3e/4+2e/64=25e/32$.
This concludes the proof of Theorem~\ref{thm:main}.

\paragraph{Search size.}
A grid point is a choice of $d$ signs and a weak composition of $D$
into $d$ parts. The bar positions in a string of $D$ marks and $d-1$
bars encode those compositions, giving the count in the theorem.
Alternatively a scan may generate signed integer tuples in $[-D,D]^d$
and filter on $\sum|c_i|=D$. Each top facet has at most $(d!)^r$
top simplices after $r$ subdivisions, by the flag formula in
Section~\ref{sec:finite-degree}. Thus $d\,2^d(d!)^r$ bounds the
number of fine vertex occurrences used in the proof. The program
returning $W(k)$ only needs the denominator-$D$ grid. Its grid size
can grow rapidly with the input modulus.

\paragraph{An object-language induction ledger.}
For definiteness, encode $(a,b)$ by $2^a(2b+1)-1$, finite lists by
their length and a nested pairing, and integers by
$0\mapsto0$, $z>0\mapsto2z$, $-z\mapsto2z-1$. Pair decoding is
bounded by the code plus one. Use a signed numerator and a positive
denominator for rationals, reducing by a bounded gcd. Operations
have defaults on malformed codes and on division by zero.
Comparisons are made by cross multiplication. Finite products of
Boolean tests express bounded universal quantifiers.
The following are the induction predicates needed in the proof.
\begin{itemize}
\item For an oriented sum, store a coefficient for every enumerated face.
After processing $j$ terms each coefficient is its signed prefix sum.
For a cancellation involution one may collect the two occurrences by
their lesser index as a bucket; the equality of the two orders of
summation follows by a rectangular double-sum calculation on prefixes.
\item At subdivision step $i$, check the coefficient equations for all
faces on the finite lists at the current stage. Use the finite face
conjunction again for the carrier proof by face dimension. For the
diameter bound, compare all vertex pairs on all enumerated simplices.
At time step $i$, compare the top-image coefficient with the initial
coefficient. All these assertions are equations or Boolean tests of PR
terms, even when $d,r$ vary.
\item For a scan of a list of length $B$, after $j$ entries a recorded
index is in range and passes the test, or the failure flag is set and
the Boolean disjunction through $j$ is false. The update either stores
the next passing index or preserves the invariant. At termination the
exhibited good grid point rules out failure.
\end{itemize}
To collect the required instances of (C), (U), (T) into a Boolean
table, keep the field hypotheses fixed and induct on the number of
instances collected, with the prefix conjunction as the induction
formula. The tail precision $j$ in Theorem~\ref{thm:main} is arbitrary
only after the witness has been computed. There is no induction on a
field contract, no choice of a modulus from continuity, and no
universal interpreter for PR programs.
